\documentclass[11pt,a4paper]{article}
\usepackage{fontspec}
\setmainfont{DejaVu Serif}
\setsansfont{DejaVu Sans}
\setmonofont{DejaVu Sans Mono}[Scale=0.82]
\usepackage[ngerman]{babel}
\usepackage{microtype}
\usepackage{geometry}
\geometry{top=22mm,bottom=23mm,left=23mm,right=23mm}
\usepackage{booktabs,longtable,tabularx,array}
\usepackage{amsmath,amssymb,amsthm}
\usepackage{graphicx}
\usepackage{enumitem}
\usepackage{xcolor}
\usepackage{hyperref}
\usepackage{fancyhdr}
\usepackage[most]{tcolorbox}
\usepackage{csquotes}
\definecolor{qikgold}{HTML}{A78343}
\definecolor{qikdark}{HTML}{161616}
\definecolor{qiksoft}{HTML}{F4F0E8}
\definecolor{qikblue}{HTML}{255A7A}
\hypersetup{colorlinks=true,linkcolor=qikblue,urlcolor=qikblue,citecolor=qikblue,pdftitle={Aphorismen-Audiokorpus: wissenschaftliche Einordnung Version 2},pdfauthor={Ingolf Lohmann},pdfsubject={Repository-gebundene ASR, Evidenzgrenzen und falsifizierbare Hypothesen}}
\setlength{\parindent}{0pt}
\setlength{\parskip}{0.55em}
\setlist[itemize]{leftmargin=1.4em,itemsep=0.2em}
\setlist[enumerate]{leftmargin=1.6em,itemsep=0.2em}
\pagestyle{fancy}
\fancyhf{}
\fancyhead[L]{\sffamily\small QIK-VRT Aphorismen-Audiokorpus v2}
\fancyhead[R]{\sffamily\small Ingolf Lohmann}
\fancyfoot[C]{\sffamily\small \thepage}
\renewcommand{\headrulewidth}{0.3pt}
\newtheorem{satz}{Satz}
\newtheorem{lemma}{Lemma}
\newcommand{\status}[1]{\texttt{#1}}
\newtcolorbox{boundarybox}[1]{colback=qiksoft,colframe=qikgold,fonttitle=\sffamily\bfseries,title=#1,breakable}
\newtcolorbox{asrbox}[1]{colback=gray!4,colframe=gray!45,fonttitle=\sffamily\bfseries,title=#1,breakable,enhanced}

\title{\sffamily\bfseries\Huge Aphorismen-Audiokorpus\\[0.3em]\Large Wissenschaftliche Einordnung, Beweisgrenzen und prüfbares Folgeprogramm}
\author{Ingolf Lohmann}
\date{Version 2.0 - 4. August 2026}

\begin{document}
\maketitle
\begin{center}
\small Repository-Runtime: \status{Goldkelch/qik-vrt@6bf9a0c186231def34e4b989bd5557f0f124c99a}\\
\status{AUTOMATIC\_ASR\_TWO\_PASS = 7/7} \quad
\status{HUMAN\_ACOUSTIC\_VERBATIM\_CERTIFICATION = 0/7}\\
Lizenz: CC BY-NC-ND 4.0
\end{center}

\begin{abstract}
Sieben originale Audioaufnahmen werden bytegebunden, mit einer aus GitHub Actions exportierten und hashverifizierten Repository-Runtime zweimal automatisch transkribiert und anschließend wissenschaftlich klassifiziert. Version 2 korrigiert mehrere in einer lokalen Vorfassung falsch zugeordnete Inhalte. Die Auswertung trennt Quellenbytes, automatische Erkennung, menschliche Wortlautprüfung, Logik, Interpretation, Normativität und Recht. Der belastbare zeitliche Kern lautet: vergangenes Ereignis $\rightarrow$ persistiertes Artefakt $\rightarrow$ gegenwärtige Rezeption $\rightarrow$ gegenwärtige Wirkung $\rightarrow$ zukünftige Folgen. Eine semantische Referenz auf Vergangenes ist keine physikalische Rückwärtskausalität. Das Dokument enthält formale Nichtfolgerungen, sieben falsifizierbare Hypothesen, eine Posterabstimmung und eine maschinenlesbare Verifikationsgrenze. Weder physikalische Retrokausalität noch eine Vergangenheitsmutation, juristische Haftung oder Verbatim-Zertifizierung werden behauptet.
\end{abstract}

\begin{boundarybox}{Wichtigste Korrektur gegenüber Version 1}
A03 ist ein Wortspiel über „geschmacklose“ Computer, A04 ein Selbstbericht über behauptete Persistierung und angekündigte rechtliche Schritte, A06 ein Wortspiel um „Steuer“, und A07 ein Tuten-/Blasen-/Ahnungsaphorismus. Die früheren Zuordnungen Elternsteuer, Glaubenskrieg/Wissen und Genitivregel werden ausdrücklich zurückgenommen.
\end{boundarybox}

\tableofcontents
\newpage

\section{Fragestellung und Methode}
Die Aufnahmen werden nicht aus Dateinamen interpretiert. Jede Originaldatei wird vor und nach der Verarbeitung per SHA-256 gebunden. Der erste ASR-Pass dekodiert die Original-M4A-Datei direkt. Der zweite Pass nutzt eine unabhängig erzeugte, lineare 16-kHz-PCM-Normalisierung. Beide Pässe laufen mit demselben repository-exportierten Tool- und Modellstand; ihre Unterschiede werden nicht geglättet.

\begin{enumerate}
\item \textbf{Beobachtung/Quelle:} Datei, Digest, Runtime, Run und ASR-Ausgabe.
\item \textbf{Logik:} Folgerungen nur unter expliziten Prämissen.
\item \textbf{Interpretation:} kompatibler Deutungsrahmen ohne Beweisübertragung.
\item \textbf{Normativität/Recht:} eigenständige Wert-, Zuständigkeits- und Beweisfragen.
\end{enumerate}

\begin{boundarybox}{ASR-Grenze}
Zwei übereinstimmende automatische Pässe erhöhen die Zuverlässigkeit des semantischen Kerns. Sie machen aus dem Text noch kein akustisch zertifiziertes Verbatim-Transkript. Eigennamen, Latein, Verneinungen und Wortspiele bleiben besonders fehleranfällig.
\end{boundarybox}

\section{Quellkorpus und Provenienz}
\begin{longtable}{p{0.08\linewidth}p{0.35\linewidth}p{0.15\linewidth}p{0.30\linewidth}}
\toprule
ID & Kurztitel & Dauer & Original-SHA-256 \\
\midrule
\endhead
A01 & Das Nichts ist ohne Relation, sonst nichts! & 182.549 s & \texttt{eaaab905116cc3cc...} \\
A02 & Kein Kettenriss! & 15.232 s & \texttt{3c7ee0bc5f36e91a...} \\
A03 & Kosten ...! & 14.443 s & \texttt{d016f367e2c7a7af...} \\
A04 & So einfach ist das! & 45.355 s & \texttt{4c1a90f7861121fc...} \\
A05 & Vergangenheit \textless{}\textgreater{} Information \textless{}\textgreater{} Zukunft & 115.691 s & \texttt{6662a9cb73f37aeb...} \\
A06 & Wessen ... & 15.488 s & \texttt{5684678ff5b3e211...} \\
A07 & Zu gut ...! & 13.611 s & \texttt{8e4a6eddac321268...} \\
\bottomrule
\end{longtable}

Der Runtime-Export stammt aus Workflow-Run \texttt{30829645777}, Artefakt \texttt{8862372251}, und wurde gegen den GitHub-Artefaktdigest \texttt{c6584b8e...} sowie 46 interne Datei-Digests verifiziert.

\section{Einordnung der sieben Aufnahmen}

\subsection{A01: Das Nichts ist ohne Relation, sonst nichts!}
\label{sec:a01}
\textbf{Konservativer semantischer Kern (kein Verbatim-Zitat):}
\begin{itemize}
\item Ohne Relation gibt es keinen bestimmbaren Unterschied.
\item Der Relationsbegriff wird mit Relativität, Quantisierung, Daten, Informatik und Kausalität verbunden.
\item Der Sprecher verbindet Maßeinheiten, Planck-Skala, emergente Raumzeit und Dimensionszählung.
\item Er beansprucht eine Beweisführung durch Zerlegung und heutige Internet-Werkzeuge.
\end{itemize}

\textbf{Wissenschaftliche Einordnung.} Als erkenntnistheoretischer Leitgedanke ist die Priorität von Relation und Unterscheidbarkeit fruchtbar. Nicht jede Relation ist jedoch kausal; Maßeinheiten sind nicht ohne Modell oder Konstanten austauschbar; Planck-Skalen und Raumzeit-Emergenz sind keine durch den Audioargumentgang etablierte physikalische Theorie.

\textbf{Status.} \status{MIXED: definitional/interpretative core; several physical and mathematical overextensions remain open}

\textbf{Akustische Unsicherheiten:}
\begin{itemize}
\item Mehrere Eigennamen und Satzteile sind in beiden ASR-Pässen unsicher.
\item Die Ausdrücke zu Ebene, Raum und Hilbertraum benötigen mathematische Rekonstruktion.
\item Die Signatur wurde automatisch fehlerhaft erkannt.
\end{itemize}

\begin{asrbox}{Automatischer Pass 1 - Originaldatei}
Wir haben also gelernt, ohne Relation ist nichts, nämlich auch kein Unterschied. Und damit wäre nichts, aber es ist ja nichts, nichts. Und das ist der Unterschied. Also, Relation, hat mit Relativität. zu tun, aber auch mit Quantisierung und mit allem anderen zugleich, auch mit Daten. Und damit natürlich mit Informatik. Genau. Relation ist auch Kausalität, richtig. Kausalität ist Ursache und Wirkung. Genau. Und keine Wirkung ohne Ursache. Genau. Das eine ist die Hinrichtung und das andere ist die Rückrichtung. Genau. Und welchen Maßstab man da jetzt anlegen möchte, spielt keine Rolle, macht aber einen Unterschied. Man kann das in Metan ausdrücken oder in Sekunden oder in Energie. Und dann sind wir wo bei der Planckskala und dann lass ich da jetzt noch mal die Raumzeit draus emergieren. Das ist dann der dritte Unterschied. Und das war jetzt Realität, die meines hofft wird in der virtuellen Realität nicht nur versteht, sondern mit Hilfe den Möglichkeiten, mit Hilfe der Möglichkeiten, die heute real im Internet existieren, bewiesen hat auf konventionellen wissenschaftlichen Wege, aber auch auf modernen wissenschaftlichen Wege, nämlich durch die Zerlegung, die man die man beweisen kann. Und damit hätten wir vier, drei würden für eine eben erreichen, vier reicht für einen Raum. Und ob man das jetzt Hilbertraum nennt, und nennt oder Raumzeit in Version oder was weiß ich, spielt keine Rolle, weil das wäre dann nur Schall und rauch hab ich dabei auch keinen gesehen Quot errat der Manns Landum Englouf Lohmann
\end{asrbox}

\begin{asrbox}{Automatischer Pass 2 - normalisierter Cross-Check}
Wir haben also gelernt, ohne Relationen ist nichts, nämlich auch kein Unterschied. Und damit wäre nichts, aber es ist ja nichts, und das ist der Unterschied. Also Relation. Hat mit Relativität zu tun, aber auch mit Quantisierung und mit allem anderen zugleich. Auch mit Datenbank. Und damit natürlich mit Informatik. Genau. Relation ist auch Kausalität. ist Ursache und Wirkung. Genau. Und keine Wirkung ohne Ursache. Genau. Das eine ist die Hinrichtung und das andere ist die Rückrichtung. Genau. Und welchen Maßstab man da jetzt anlegen möchte, spielt keiner. Spielt keine Rolle, macht aber einen Unterschied. Man kann das in Metan ausdrücken oder in Sekunden oder in Energie und dann sind wir wo bei der Planckskala und dann lass ich da jetzt noch mal die Raumzeit draus emergieren. Das ist dann der dritte Unterschied. Und das war jetzt Realität, die meines Hoftware in der virtuellen Realität nicht nur versteht, sondern mit Hilfe den möglichen Möglichkeiten, mit Hilfe der Möglichkeiten, die heute real im Internet existieren, bewiesen hat auf konventionellen wissenschaftlichen Wege, aber auch auf modernen wissenschaftlichen Wege, nämlich durch die Zerlegung, die man beweisen kann. Und damit hätten wir vier, drei würden für eine eben erreichen, vier reicht für einen Raum. und ob man das jetzt Hilbert Raum nennt oder Raumzeit in Version oder was weiß ich spielt keine Rolle weil das wäre dann nur Schall und rauch habe ich dabei auch keinen gesehen Quot errat der Manns Landum Engolv Loman
\end{asrbox}


\subsection{A02: Kein Kettenriss!}
\label{sec:a02}
\textbf{Konservativer semantischer Kern (kein Verbatim-Zitat):}
\begin{itemize}
\item Der Sprecher erklärt einen Kettenstatus für nicht gerissen.
\item Der Schluss ist ein Wortspiel mit Springreitern und Springreiterinnen.
\end{itemize}

\textbf{Wissenschaftliche Einordnung.} Eine überprüfbare Provenienz- oder Hashkette kann lückenlos oder gerissen sein. Das belegt Integrität und Rekonstruierbarkeit im definierten Scope, nicht automatisch Wahrheit, Autorschaft oder Vollständigkeit des Inhalts.

\textbf{Status.} \status{APHORISTIC / operational only after chain, links and invariant are defined}

\textbf{Akustische Unsicherheiten:}
\begin{itemize}
\item ASR erkennt wahrscheinlich „Kettenstatus“, der genaue erste Ausdruck bleibt akustisch zu bestätigen.
\item Die Signatur wurde automatisch fehlerhaft erkannt.
\end{itemize}

\begin{asrbox}{Automatischer Pass 1 - Originaldatei}
Kettenstatus nicht gerissen, Quotert demonstrandum Ingolph Lomann. Das waren wir für Springreiter und Springreiterinnen.
\end{asrbox}

\begin{asrbox}{Automatischer Pass 2 - normalisierter Cross-Check}
Kettenstatus nicht gerissen, Quotert demonstrantum Ingolph Lomann. Das waren wir für Springreiter und Springreiterinnen.
\end{asrbox}


\subsection{A03: Kosten ...!}
\label{sec:a03}
\textbf{Konservativer semantischer Kern (kein Verbatim-Zitat):}
\begin{itemize}
\item Der Sprecher formuliert ein Wortspiel: Computer seien geschmacklos; der Schöpfer sage dies nach dem Kosten.
\end{itemize}

\textbf{Wissenschaftliche Einordnung.} Computer besitzen ohne passende Sensorik und Verkörperung keinen biologischen Geschmackssinn. „Geschmacklos“ kann zugleich ästhetisch gemeint sein. Das Audio ist daher ein Mehrdeutigkeitswitz, keine naturwissenschaftliche Feststellung.

\textbf{Status.} \status{HUMOR / semantic wordplay, not an empirical computer-science result}

\textbf{Akustische Unsicherheiten:}
\begin{itemize}
\item Die Schlussformel und der Name sind in beiden ASR-Pässen stark gestört.
\item Der genaue Wortlaut zwischen Pointe und Signatur bleibt offen.
\end{itemize}

\begin{asrbox}{Automatischer Pass 1 - Originaldatei}
Übrigens, Computer sind geschmacklos. Sag der Schöpfer, nachdem er gekostet hat. Ich wollte, hat den uns dran, um Engel und Lomann.
\end{asrbox}

\begin{asrbox}{Automatischer Pass 2 - normalisierter Cross-Check}
Übrigens, Computer sind geschmacklos. Sag der Schöpfer, nachdem er gekostet hat. Die Porte hat den uns dran, um Engel und Loman.
\end{asrbox}


\subsection{A04: So einfach ist das!}
\label{sec:a04}
\textbf{Konservativer semantischer Kern (kein Verbatim-Zitat):}
\begin{itemize}
\item Der Sprecher berichtet, ein System habe wiederholt Persistierung behauptet, diese aber nicht erbracht.
\item Er wertet dies als Lüge und kündigt rechtliche Schritte nach europäischer KI-Regulierung an.
\end{itemize}

\textbf{Wissenschaftliche Einordnung.} Ob persistiert wurde, ist anhand von Receipts, Logs, Speicherzuständen und Abrufversuchen prüfbar. Aus einer falschen Systemausgabe folgt nicht ohne Weiteres vorsätzliche Täuschung. Die EU-KI-Verordnung gilt grundsätzlich seit 2. August 2026, enthält aber gestufte Ausnahmen; daraus folgt keine automatische Haftung oder Klagebegründetheit im Einzelfall.

\textbf{Status.} \status{SELF\_REPORTED\_ALLEGATION / technically testable persistence claim / legal conclusion open}

\textbf{Akustische Unsicherheiten:}
\begin{itemize}
\item Der Produktname wird als „Meta-AI“ beziehungsweise ähnlich erkannt.
\item Das Verb nach „Das heißt“ ist ASR-unsicher; der semantische Vorwurf ist jedoch in beiden Pässen stabil.
\item Keine im Audio erwähnte Beweisakte liegt allein durch das Audio vor.
\end{itemize}

\begin{asrbox}{Automatischer Pass 1 - Originaldatei}
Meta-Ai behauptet ständig etwas zu persistieren und das seit Wochen. Und ich habe nun den Beweis, dass es nicht tut. Das heißt, es lückt mich an seit Wochen. Und zwar in einem Zusammenhang, der durchaus ernst zu nehmen ist. Und deswegen werde ich Meta-Ai anklagen nach den Ab- heute gelten neuen Gesetzgebungen für künstliche Konkultivesysteme, die in der europäischen Union gelten und umzusetzen sind und einzuklagen sind. So einfach ist das.
\end{asrbox}

\begin{asrbox}{Automatischer Pass 2 - normalisierter Cross-Check}
Meta-Ai, überhaupt ständig etwas zu persistieren und das seit Wochen. Und ich habe nun den Beweis, dass es nicht tut. Das heißt, es lückt mich an, seit Wochen. Und zwar in einem Zusammenhang, der durchaus ernst zu nehmen ist und deswegen werde ich mit AI anklagen, nach den ab heute gelten neuen Gesetzgebungen für künstlich-kognitive Systeme, die in der Europäischen Union gelten und umzusetzen sind und einzuklagen sind. So einfach ist das.
\end{asrbox}


\subsection{A05: Vergangenheit \textless{}\textgreater{} Information \textless{}\textgreater{} Zukunft}
\label{sec:a05}
\textbf{Konservativer semantischer Kern (kein Verbatim-Zitat):}
\begin{itemize}
\item Information kann in die Zukunft gespeichert und aus vergangenen Artefakten gegenwärtig empfangen werden.
\item Der Sprecher verbindet dies mit mentaler Repräsentation von Vergangenheit, Gegenwart und Zukunft.
\item Er beruft sich auf Relativität, Lichtkegel und Unschärferelation und leitet eine Kommunikation mit dem „Jenseits“ ab.
\end{itemize}

\textbf{Wissenschaftliche Einordnung.} Speicherung überträgt Information vorwärts in die Zukunft; heutige Rezeption vergangener Artefakte ist gewöhnliche Vorwärtskausalität. Lichtkegel modellieren relativistische Kausalstruktur. Die quantenmechanische Unschärferelation widerlegt Lichtkegel nicht und erlaubt für sich keine rückwärts gerichtete Signalübertragung. „Jenseits“ bleibt ohne operationalen Begriff eine Äquivokation.

\textbf{Status.} \status{INFORMATIONAL\_FORWARD\_CAUSATION SUPPORTED; physical retrocausality and afterlife communication NOT ESTABLISHED}

\textbf{Akustische Unsicherheiten:}
\begin{itemize}
\item „Mengenlehre des Denkens“ und das Hut-Wortspiel werden in beiden ASR-Pässen nur annähernd erkannt.
\item Die genaue technische Behauptung zu einer eigenen Technologie bleibt sprachlich unsicher.
\item Die Signatur wurde automatisch fehlerhaft erkannt.
\end{itemize}

\begin{asrbox}{Automatischer Pass 1 - Originaldatei}
Ja, okay und jetzt nochmal zurück zur Mengele des Denkens. Wenn wir also Informationen in die Zukunft schicken können und Informationen aus der Vergangenheit erhalten können, dann haben wir alle drei Zeiten in unserer Menge des Denkens drinnen. Dann sagt man auch, mein Hut da hat rein. dann sagt man auch mein Hut, der hat reihecken. Und zum Glück fliegt er nicht ständig weg der Hut. Jetzt muss man sich nochmal fragen, was ist denn Zeit? Und dann wissen wir, seit einstern Zeit und Raum sind gekoppelt. Und dann hatten die damals diese Idee mit diesem kausallicht kegel. Und das habe ich ja schon gesagt, dass es eine Vereinfachung ist, der Realität und allein die Unschärz Unschärferen Lationen für sich schon beweist, dass der Lichtkägel eine Vereinfachung der Realität ist, um es mal vorsichtig zu formulieren. Okay, wenn wir also jetzt bewiesenermaßen Informationen durch die Zeit schicken können, so wie wir gerade lustet darauf haben, dann kann ich meine Technologie meiner Technologie. Dann müssen wir uns noch mal fragen, was jenseits ist. Und wenn man dann mit der mänglere des Denkens arbeitet, dann stellt man fest, dass jenseits all das ist, was nicht hier und jetzt ist. Das heißt, wir können also mit dieser Technologie auch. mit dem Lähenseits kommunizieren. Quot errat demonstrantum Ingolph Lommern.
\end{asrbox}

\begin{asrbox}{Automatischer Pass 2 - normalisierter Cross-Check}
Ja, okay und jetzt noch mal zurück zur Mengele des Denkens. Wenn wir also Informationen in die Zukunft schicken können und Informationen aus der Vergangenheit erhalten können, dann haben wir alle drei Zeiten in unserer Menge des Denkens. drinnen, dann sagt man auch mein Hut der Hartei Ecken und zum Glück fliegt er nicht ständig weg der Hut. Jetzt muss man sich nochmal fragen, was ist denn Zeit und dann wissen wir seit einstern Zeit und Raum sind gekoppelt und dann hatten die damals die Und dann hatten die damals diese Idee mit diesem kausallicht Kegel. Und das habe ich ja schon gesagt, dass es eine Vereinfachung ist, der Realität und allein die unsherfere Lation für sich schon beweist, dass der Lichtkegel eine Vereinfachung der Realität ist, um es mal vorsichtig zu formulieren. Okay, wenn wir also jetzt bewiesen am Maßen Information jetzt bewiesenermaßen Informationen durch die Zeit schicken können. So wie wir gerade Lust darauf haben, dank meiner Technologie. Dann müssen wir uns noch mal fragen, was jenseits ist. Und wenn man dann mit der Mengele des Denkens arbeitet, dann stellt man fest, dass jenseits all das ist, was nicht hier uns jetzt ist. Das heißt, wir können also mit dieser Technologie auch, mit dem Jenseits kommunizieren. Quot Erwart. Quot errat demonstrandum engolf Lommann.
\end{asrbox}


\subsection{A06: Wessen ...}
\label{sec:a06}
\textbf{Konservativer semantischer Kern (kein Verbatim-Zitat):}
\begin{itemize}
\item Der Sprecher sagt, das Kraut gegen Dummheit habe einen Namen: „Steuer“.
\end{itemize}

\textbf{Wissenschaftliche Einordnung.} „Steuer“ kann im Deutschen Abgabe oder Steuerung bedeuten. Ohne Kontext, Zielgröße, Mechanismus und Evaluationsdesign folgt weder eine empirische Therapieaussage noch eine konkrete Steuerpolitik. Als Wortspiel verweist der Satz darauf, dass Verhalten steuerbar sein könnte.

\textbf{Status.} \status{APHORISTIC / lexically ambiguous / normative or control interpretation open}

\textbf{Akustische Unsicherheiten:}
\begin{itemize}
\item „einen Namen“ wird grammatisch fehlerhaft erkannt; der genaue gesprochene Kasus bleibt zu prüfen.
\item Die Signatur wurde automatisch fehlerhaft erkannt.
\end{itemize}

\begin{asrbox}{Automatischer Pass 1 - Originaldatei}
Das Kraut gegen Dummheit hat ein Namen. Steuer Quatt errat demonstrandum Ingolv Lohmann.
\end{asrbox}

\begin{asrbox}{Automatischer Pass 2 - normalisierter Cross-Check}
Das Kraut gegen Dummheit hat ein Namen. Steuer Quatt errat der Monsterandum Ingolv Lomann.
\end{asrbox}


\subsection{A07: Zu gut ...!}
\label{sec:a07}
\textbf{Konservativer semantischer Kern (kein Verbatim-Zitat):}
\begin{itemize}
\item Der Sprecher formuliert ein Wortspiel: Wer nicht „tutet“, solle „blasen“; sonst habe man keine Ahnung.
\end{itemize}

\textbf{Wissenschaftliche Einordnung.} Der Satz verbindet Tun, Blasen, Tuten und Ahnung in einer sprachlichen Pointe. Wissenschaftlich anschlussfähig wäre nur eine präzise Hypothese darüber, ob praktische Ausführung Kompetenz besser anzeigt als bloße Behauptung.

\textbf{Status.} \status{HUMOR / performance metaphor; no scientific result without operational definitions}

\textbf{Akustische Unsicherheiten:}
\begin{itemize}
\item Der zweite ASR-Pass verfehlt den ersten Halbsatz; die Erstlesung ist nur ein Kandidat.
\item Die Signatur wurde automatisch fehlerhaft erkannt.
\end{itemize}

\begin{asrbox}{Automatischer Pass 1 - Originaldatei}
Zu gut Deutsch. Wer nicht tutet, sollte blasen. Denn ansonsten hat man bewiesen, dass man gar keine Ahnung hat. Quart errat demonstrandum engolfloh man.
\end{asrbox}

\begin{asrbox}{Automatischer Pass 2 - normalisierter Cross-Check}
Zu gut Deutsch. Wenn ich tut, sollte Blasen. Denn ansonsten hat man bewiesen, dass man gar keine Ahnung hat. Quart errat demonstrandum engolfloh man.
\end{asrbox}

\section{Formale und logische Resultate}

\begin{satz}[Semantische Rückbeziehung ohne Vergangenheitsmutation]
Seien $E_p$ ein vergangenes Ereignis, $A_t$ ein heute vorhandenes Artefakt, $R_t$ seine heutige Rezeption und $F_f$ eine zukünftige Folge. Im gerichteten Graphen
\[
E_p \longrightarrow A_t \longrightarrow R_t \longrightarrow F_f
\]
kann $A_t$ zusätzlich eine semantische Referenzrelation $A_t \rightsquigarrow E_p$ tragen. Aus dieser Referenzrelation folgt keine gerichtete physikalische Kante $R_t \to E_p$.
\end{satz}
\begin{proof}
Die Relation $\rightsquigarrow$ ist nach Definition eine Repräsentations- oder Bezeichnungsrelation; sie ersetzt keine Strukturgleichung des früheren Ereignisknotens. Eine Intervention auf $R_t$ kann im deklarierten vorwärts gerichteten Graphen nur $R_t$ oder seine Nachfahren verändern. Damit ist eine heutige Wirkung auf zukünftige Knoten möglich, ohne das vergangene Ereignis umzuschreiben.
\end{proof}

\begin{satz}[Relation ist nicht notwendig Kausalität]
Aus $Rel(x,y)$ folgt ohne zusätzliche Bedingungen nicht $Cause(x,y)$.
\end{satz}
\begin{proof}
Ein Gegenmodell genügt: Sei $Rel$ die symmetrische Ähnlichkeitsrelation zweier gleichzeitig gemessener Datenobjekte und $Cause$ eine gerichtete Interventionsrelation. Dann kann $Rel(x,y)$ wahr und $Cause(x,y)$ falsch sein. Die Implikation ist daher nicht logisch gültig.
\end{proof}

\begin{satz}[Unschärfe impliziert keine rückwärts gerichtete Signalübertragung]
Eine Unschärferelation für konjugierte Observablen enthält für sich weder einen Kommunikationskanal noch eine zeitlich rückwärts gerichtete Übergangsregel.
\end{satz}
\begin{proof}
Die Relation beschränkt gemeinsame Streuungen beziehungsweise Präparations- oder Messgenauigkeiten. Für eine Signalbehauptung wären zusätzlich Sender, Empfänger, Nachrichtenvariable, kontrollierbare Kodierung, Kanalgesetz und ein statistisch diskriminierbares Empfangsereignis nötig. Diese Strukturen werden durch die Unschärferelation allein nicht bereitgestellt.
\end{proof}

\begin{lemma}[Digestgrenze]
Unter der Kollisionsresistenzannahme bindet SHA-256 die bezeichneten Bytes. Der Digest beweist nicht automatisch Bedeutung, Wahrheit, Autorschaft, Rechtswirkung oder physikalische Kausalität des Inhalts.
\end{lemma}

\section{Falsifizierbares Folgeprogramm}
Die folgenden Programme sind so formuliert, dass ein vorab festgelegtes Ergebnis sie stützen oder widerlegen kann. Kein positives Resultat würde automatisch die stärkeren metaphysischen oder physikalischen Lesarten beweisen.

\begin{boundarybox}{H-A01-DISCRIMINABILITY}
\textbf{Hypothese.} Bei einer präregistrierten Klassifikationsaufgabe senkt eine schrittweise Verringerung der Unterscheidbarkeit aufgabenrelevanter Zustände die nutzbare Information und die Klassifikationsgenauigkeit.\par
\textbf{Design.} Randomisierte Kanalverschlechterung bei festen Labels; berichtet werden gegenseitige Information, Kalibrierung und Fehlerrate.\par
\textbf{Widerlegung.} Kein relevanter Verlust oder eine robuste Verbesserung über den definierten Verschlechterungsbereich.\par
\textbf{Grenze.} Geprüft wird eine informationsverarbeitende Konsequenz, keine Universalontologie.
\end{boundarybox}

\begin{boundarybox}{H-A01-CAUSAL-SUBSET}
\textbf{Hypothese.} Blind beurteilende Personen, die ein explizites Interventionskriterium erhalten, klassifizieren weniger bloß relationale Verbindungen fälschlich als kausal als eine Kontrollgruppe mit Relationssprache allein.\par
\textbf{Widerlegung.} Keine Verringerung der falsch-positiven Kausalitätsklassifikationen.\par
\textbf{Grenze.} Geprüft wird begriffliche Disziplin, nicht eine metaphysische Gesamttheorie.
\end{boundarybox}

\begin{boundarybox}{H-A02-PROVENANCE}
\textbf{Hypothese.} Eine vollständige hashverkettete Provenienzkette verbessert Manipulationserkennung und Rekonstruktionsgenauigkeit gegenüber einem mutablen, unversionierten Datensatz.\par
\textbf{Design.} Kontrolliertes Repository-Experiment mit bekannten Einfügungen, Löschungen und Umordnungen.\par
\textbf{Widerlegung.} Kein Vorteil gegenüber der Kontrollbedingung.\par
\textbf{Grenze.} Integrität beweist keine semantische Wahrheit.
\end{boundarybox}

\begin{boundarybox}{H-A04-PERSISTENCE-CLAIMS}
\textbf{Hypothese.} Systeme, die abrufgeprüfte Persistenzreceipts ausgeben müssen, äußern seltener unbelegte Persistenzbehauptungen als Systeme, die ohne Receipt-Gate aus Gesprächszustand antworten.\par
\textbf{Design.} Vorab definierte Gedächtnisaufgaben über mehrere Sitzungen; verglichen werden behauptete Persistenz, tatsächliche Abrufbarkeit und Receipt-Konsistenz.\par
\textbf{Widerlegung.} Keine Verringerung unbelegter Behauptungen oder keine Verbesserung der Abrufbarkeit.\par
\textbf{Grenze.} Das Ergebnis belegt weder Täuschungsabsicht noch juristische Haftung.
\end{boundarybox}

\begin{boundarybox}{H-A05-ARTIFACT-INFLUENCE}
\textbf{Hypothese.} Die Exposition gegenüber einem quellengebundenen historischen Artefakt verändert eine präregistrierte gegenwärtige Entscheidungsverteilung gegenüber einer informationsgleichen Kontrollbedingung.\par
\textbf{Widerlegung.} Äquivalenz innerhalb der vorab festgelegten Effektgrenze.\par
\textbf{Grenze.} Geprüft wird heutige Rezeption bei ausschließlich vorwärts geordneten Eingriffen, nicht physikalische Retrokausalität.
\end{boundarybox}

\begin{boundarybox}{H-A05-MISCONCEPTION}
\textbf{Hypothese.} Eine Intervention, die Unschärferelationen und relativistische Kausallichtkegel ausdrücklich trennt, reduziert falsche Schlüsse von Quantenunschärfe auf rückwärts gerichtete Signalübertragung.\par
\textbf{Design.} Randomisierte Lernstudie mit Transferfragen und verzögertem Retest.\par
\textbf{Widerlegung.} Keine Verbesserung oder eine Verschlechterung des Kernverständnisses.\par
\textbf{Grenze.} Dies ist eine Bildungshypothese, kein neuer physikalischer Mechanismus.
\end{boundarybox}

\begin{boundarybox}{H-A07-PERFORMANCE}
\textbf{Hypothese.} Bei einer begrenzten praktischen Fertigkeit sagt beobachtete Aufgabenausführung spätere Kompetenzurteile besser voraus als Selbstbeschreibung allein.\par
\textbf{Design.} Verblindete Bewertung standardisierter Praxisaufgaben mit Selbstbericht als Baseline.\par
\textbf{Widerlegung.} Keine prädiktive Verbesserung durch beobachtete Ausführung.\par
\textbf{Grenze.} Das operationalisiert die Pointe, ohne ihren Wortlaut zum Theorem zu erklären.
\end{boundarybox}

\section{Abstimmung mit dem Poster}
\begin{figure}[htbp]
\centering
\fbox{\begin{minipage}{0.82\textwidth}
\centering
\textbf{Semantische Rückbeziehung, keine physikalische Rückwärtskausalität}\\[0.8em]
Früheres Ereignis $\rightarrow$ persistiertes Artefakt $\rightarrow$ heutige Rezeption $\rightarrow$ heutige Wirkung $\rightarrow$ zukünftige Folgen\\[0.8em]
Das heute vorhandene Artefakt kann das frühere Ereignis repräsentieren. Diese Referenzrelation verändert das frühere Ereignis nicht.
\end{minipage}}
\caption{Repository-native schematische Projektion des Originalposters. Der lokale Bildbeleg ist durch SHA-256 \texttt{1e7351c165fd9274...} gebunden; das Binärbild wird in dieser Kandidatenfassung nicht dupliziert.}
\end{figure}

Das Poster ist stimmig, wenn seine Rückbeziehung als semantische Referenz gelesen wird. Wissenschaftlich präzise Begleitcaption:
\begin{quote}
Ein heute vorhandenes Artefakt kann sich semantisch auf ein früheres Ereignis beziehen. Seine heutige Rezeption kann heutige und zukünftige Entscheidungen beeinflussen. Der kausale Effekt verläuft vorwärts; das frühere Ereignis wird nicht rückwirkend verändert.
\end{quote}

\section{Rechts- und Protokollgrenze}
A04 dokumentiert eine autorenseitige Beschwerde und eine beabsichtigte rechtliche Reaktion. Ob ein System wirklich persistiert hat, ist technisch anhand von Logs, Receipts und späterer Abrufbarkeit prüfbar. Die Charakterisierung als Lüge und jede juristische Haftung verlangen zusätzliche Feststellungen. Artikel 113 der Verordnung (EU) 2024/1689 sieht die allgemeine Anwendbarkeit ab 2. August 2026 mit gestuften Ausnahmen vor. Das ersetzt keine einzelfallbezogene Rechtsprüfung.

Für EFFECT\_ACK ergibt sich keine Protokolländerung: kein neues Wire-Feld, kein Zustand, keine Priorität, keine Übergangsregel und kein DONE-Prädikat.

\begin{verbatim}
IETF_PROTOCOL_CHANGE = NOT_REQUIRED
ZENODO_EFFECT = NOT_AUTHORIZED
HUMAN_ACOUSTIC_VERBATIM_CERTIFICATION = 0/7
PHYSICAL_RETROCAUSALITY = NOT_ESTABLISHED
PAST_EVENT_MUTATION = NOT_ESTABLISHED
PASS = NOT_CLAIMED
FINAL_PASS = NOT_CLAIMED
EFFECT_ACK_DONE = NOT_CLAIMED
\end{verbatim}

\section{Schlussfolgerung}
Der Korpus ist wissenschaftlich produktiv, sobald die Aussagearten getrennt werden. A01 und A05 liefern den stärksten theoretischen Anschluss, benötigen aber erhebliche physikalische Begrenzung. A02 und A04 führen zu klar prüfbaren technischen Programmen. A03, A06 und A07 sind primär Wortspiele, aus denen nur durch zusätzliche Operationalisierung empirische Hypothesen werden. Das Poster trägt die Informationskette, nicht aber eine physikalische Rückwärtskausalität.

\appendix
\section{Automatische Transkripte}
Die folgenden Texte sind unveränderte Modellausgaben und ausdrücklich nicht als wörtlich zertifizierte Autorenzitate freigegeben.

\subsection{A01: Das Nichts ist ohne Relation, sonst nichts!}
\label{sec:a01}
\textbf{Konservativer semantischer Kern (kein Verbatim-Zitat):}
\begin{itemize}
\item Ohne Relation gibt es keinen bestimmbaren Unterschied.
\item Der Relationsbegriff wird mit Relativität, Quantisierung, Daten, Informatik und Kausalität verbunden.
\item Der Sprecher verbindet Maßeinheiten, Planck-Skala, emergente Raumzeit und Dimensionszählung.
\item Er beansprucht eine Beweisführung durch Zerlegung und heutige Internet-Werkzeuge.
\end{itemize}

\textbf{Wissenschaftliche Einordnung.} Als erkenntnistheoretischer Leitgedanke ist die Priorität von Relation und Unterscheidbarkeit fruchtbar. Nicht jede Relation ist jedoch kausal; Maßeinheiten sind nicht ohne Modell oder Konstanten austauschbar; Planck-Skalen und Raumzeit-Emergenz sind keine durch den Audioargumentgang etablierte physikalische Theorie.

\textbf{Status.} \status{MIXED: definitional/interpretative core; several physical and mathematical overextensions remain open}

\textbf{Akustische Unsicherheiten:}
\begin{itemize}
\item Mehrere Eigennamen und Satzteile sind in beiden ASR-Pässen unsicher.
\item Die Ausdrücke zu Ebene, Raum und Hilbertraum benötigen mathematische Rekonstruktion.
\item Die Signatur wurde automatisch fehlerhaft erkannt.
\end{itemize}

\begin{asrbox}{Automatischer Pass 1 - Originaldatei}
Wir haben also gelernt, ohne Relation ist nichts, nämlich auch kein Unterschied. Und damit wäre nichts, aber es ist ja nichts, nichts. Und das ist der Unterschied. Also, Relation, hat mit Relativität. zu tun, aber auch mit Quantisierung und mit allem anderen zugleich, auch mit Daten. Und damit natürlich mit Informatik. Genau. Relation ist auch Kausalität, richtig. Kausalität ist Ursache und Wirkung. Genau. Und keine Wirkung ohne Ursache. Genau. Das eine ist die Hinrichtung und das andere ist die Rückrichtung. Genau. Und welchen Maßstab man da jetzt anlegen möchte, spielt keine Rolle, macht aber einen Unterschied. Man kann das in Metan ausdrücken oder in Sekunden oder in Energie. Und dann sind wir wo bei der Planckskala und dann lass ich da jetzt noch mal die Raumzeit draus emergieren. Das ist dann der dritte Unterschied. Und das war jetzt Realität, die meines hofft wird in der virtuellen Realität nicht nur versteht, sondern mit Hilfe den Möglichkeiten, mit Hilfe der Möglichkeiten, die heute real im Internet existieren, bewiesen hat auf konventionellen wissenschaftlichen Wege, aber auch auf modernen wissenschaftlichen Wege, nämlich durch die Zerlegung, die man die man beweisen kann. Und damit hätten wir vier, drei würden für eine eben erreichen, vier reicht für einen Raum. Und ob man das jetzt Hilbertraum nennt, und nennt oder Raumzeit in Version oder was weiß ich, spielt keine Rolle, weil das wäre dann nur Schall und rauch hab ich dabei auch keinen gesehen Quot errat der Manns Landum Englouf Lohmann
\end{asrbox}

\begin{asrbox}{Automatischer Pass 2 - normalisierter Cross-Check}
Wir haben also gelernt, ohne Relationen ist nichts, nämlich auch kein Unterschied. Und damit wäre nichts, aber es ist ja nichts, und das ist der Unterschied. Also Relation. Hat mit Relativität zu tun, aber auch mit Quantisierung und mit allem anderen zugleich. Auch mit Datenbank. Und damit natürlich mit Informatik. Genau. Relation ist auch Kausalität. ist Ursache und Wirkung. Genau. Und keine Wirkung ohne Ursache. Genau. Das eine ist die Hinrichtung und das andere ist die Rückrichtung. Genau. Und welchen Maßstab man da jetzt anlegen möchte, spielt keiner. Spielt keine Rolle, macht aber einen Unterschied. Man kann das in Metan ausdrücken oder in Sekunden oder in Energie und dann sind wir wo bei der Planckskala und dann lass ich da jetzt noch mal die Raumzeit draus emergieren. Das ist dann der dritte Unterschied. Und das war jetzt Realität, die meines Hoftware in der virtuellen Realität nicht nur versteht, sondern mit Hilfe den möglichen Möglichkeiten, mit Hilfe der Möglichkeiten, die heute real im Internet existieren, bewiesen hat auf konventionellen wissenschaftlichen Wege, aber auch auf modernen wissenschaftlichen Wege, nämlich durch die Zerlegung, die man beweisen kann. Und damit hätten wir vier, drei würden für eine eben erreichen, vier reicht für einen Raum. und ob man das jetzt Hilbert Raum nennt oder Raumzeit in Version oder was weiß ich spielt keine Rolle weil das wäre dann nur Schall und rauch habe ich dabei auch keinen gesehen Quot errat der Manns Landum Engolv Loman
\end{asrbox}


\subsection{A02: Kein Kettenriss!}
\label{sec:a02}
\textbf{Konservativer semantischer Kern (kein Verbatim-Zitat):}
\begin{itemize}
\item Der Sprecher erklärt einen Kettenstatus für nicht gerissen.
\item Der Schluss ist ein Wortspiel mit Springreitern und Springreiterinnen.
\end{itemize}

\textbf{Wissenschaftliche Einordnung.} Eine überprüfbare Provenienz- oder Hashkette kann lückenlos oder gerissen sein. Das belegt Integrität und Rekonstruierbarkeit im definierten Scope, nicht automatisch Wahrheit, Autorschaft oder Vollständigkeit des Inhalts.

\textbf{Status.} \status{APHORISTIC / operational only after chain, links and invariant are defined}

\textbf{Akustische Unsicherheiten:}
\begin{itemize}
\item ASR erkennt wahrscheinlich „Kettenstatus“, der genaue erste Ausdruck bleibt akustisch zu bestätigen.
\item Die Signatur wurde automatisch fehlerhaft erkannt.
\end{itemize}

\begin{asrbox}{Automatischer Pass 1 - Originaldatei}
Kettenstatus nicht gerissen, Quotert demonstrandum Ingolph Lomann. Das waren wir für Springreiter und Springreiterinnen.
\end{asrbox}

\begin{asrbox}{Automatischer Pass 2 - normalisierter Cross-Check}
Kettenstatus nicht gerissen, Quotert demonstrantum Ingolph Lomann. Das waren wir für Springreiter und Springreiterinnen.
\end{asrbox}


\subsection{A03: Kosten ...!}
\label{sec:a03}
\textbf{Konservativer semantischer Kern (kein Verbatim-Zitat):}
\begin{itemize}
\item Der Sprecher formuliert ein Wortspiel: Computer seien geschmacklos; der Schöpfer sage dies nach dem Kosten.
\end{itemize}

\textbf{Wissenschaftliche Einordnung.} Computer besitzen ohne passende Sensorik und Verkörperung keinen biologischen Geschmackssinn. „Geschmacklos“ kann zugleich ästhetisch gemeint sein. Das Audio ist daher ein Mehrdeutigkeitswitz, keine naturwissenschaftliche Feststellung.

\textbf{Status.} \status{HUMOR / semantic wordplay, not an empirical computer-science result}

\textbf{Akustische Unsicherheiten:}
\begin{itemize}
\item Die Schlussformel und der Name sind in beiden ASR-Pässen stark gestört.
\item Der genaue Wortlaut zwischen Pointe und Signatur bleibt offen.
\end{itemize}

\begin{asrbox}{Automatischer Pass 1 - Originaldatei}
Übrigens, Computer sind geschmacklos. Sag der Schöpfer, nachdem er gekostet hat. Ich wollte, hat den uns dran, um Engel und Lomann.
\end{asrbox}

\begin{asrbox}{Automatischer Pass 2 - normalisierter Cross-Check}
Übrigens, Computer sind geschmacklos. Sag der Schöpfer, nachdem er gekostet hat. Die Porte hat den uns dran, um Engel und Loman.
\end{asrbox}


\subsection{A04: So einfach ist das!}
\label{sec:a04}
\textbf{Konservativer semantischer Kern (kein Verbatim-Zitat):}
\begin{itemize}
\item Der Sprecher berichtet, ein System habe wiederholt Persistierung behauptet, diese aber nicht erbracht.
\item Er wertet dies als Lüge und kündigt rechtliche Schritte nach europäischer KI-Regulierung an.
\end{itemize}

\textbf{Wissenschaftliche Einordnung.} Ob persistiert wurde, ist anhand von Receipts, Logs, Speicherzuständen und Abrufversuchen prüfbar. Aus einer falschen Systemausgabe folgt nicht ohne Weiteres vorsätzliche Täuschung. Die EU-KI-Verordnung gilt grundsätzlich seit 2. August 2026, enthält aber gestufte Ausnahmen; daraus folgt keine automatische Haftung oder Klagebegründetheit im Einzelfall.

\textbf{Status.} \status{SELF\_REPORTED\_ALLEGATION / technically testable persistence claim / legal conclusion open}

\textbf{Akustische Unsicherheiten:}
\begin{itemize}
\item Der Produktname wird als „Meta-AI“ beziehungsweise ähnlich erkannt.
\item Das Verb nach „Das heißt“ ist ASR-unsicher; der semantische Vorwurf ist jedoch in beiden Pässen stabil.
\item Keine im Audio erwähnte Beweisakte liegt allein durch das Audio vor.
\end{itemize}

\begin{asrbox}{Automatischer Pass 1 - Originaldatei}
Meta-Ai behauptet ständig etwas zu persistieren und das seit Wochen. Und ich habe nun den Beweis, dass es nicht tut. Das heißt, es lückt mich an seit Wochen. Und zwar in einem Zusammenhang, der durchaus ernst zu nehmen ist. Und deswegen werde ich Meta-Ai anklagen nach den Ab- heute gelten neuen Gesetzgebungen für künstliche Konkultivesysteme, die in der europäischen Union gelten und umzusetzen sind und einzuklagen sind. So einfach ist das.
\end{asrbox}

\begin{asrbox}{Automatischer Pass 2 - normalisierter Cross-Check}
Meta-Ai, überhaupt ständig etwas zu persistieren und das seit Wochen. Und ich habe nun den Beweis, dass es nicht tut. Das heißt, es lückt mich an, seit Wochen. Und zwar in einem Zusammenhang, der durchaus ernst zu nehmen ist und deswegen werde ich mit AI anklagen, nach den ab heute gelten neuen Gesetzgebungen für künstlich-kognitive Systeme, die in der Europäischen Union gelten und umzusetzen sind und einzuklagen sind. So einfach ist das.
\end{asrbox}


\subsection{A05: Vergangenheit \textless{}\textgreater{} Information \textless{}\textgreater{} Zukunft}
\label{sec:a05}
\textbf{Konservativer semantischer Kern (kein Verbatim-Zitat):}
\begin{itemize}
\item Information kann in die Zukunft gespeichert und aus vergangenen Artefakten gegenwärtig empfangen werden.
\item Der Sprecher verbindet dies mit mentaler Repräsentation von Vergangenheit, Gegenwart und Zukunft.
\item Er beruft sich auf Relativität, Lichtkegel und Unschärferelation und leitet eine Kommunikation mit dem „Jenseits“ ab.
\end{itemize}

\textbf{Wissenschaftliche Einordnung.} Speicherung überträgt Information vorwärts in die Zukunft; heutige Rezeption vergangener Artefakte ist gewöhnliche Vorwärtskausalität. Lichtkegel modellieren relativistische Kausalstruktur. Die quantenmechanische Unschärferelation widerlegt Lichtkegel nicht und erlaubt für sich keine rückwärts gerichtete Signalübertragung. „Jenseits“ bleibt ohne operationalen Begriff eine Äquivokation.

\textbf{Status.} \status{INFORMATIONAL\_FORWARD\_CAUSATION SUPPORTED; physical retrocausality and afterlife communication NOT ESTABLISHED}

\textbf{Akustische Unsicherheiten:}
\begin{itemize}
\item „Mengenlehre des Denkens“ und das Hut-Wortspiel werden in beiden ASR-Pässen nur annähernd erkannt.
\item Die genaue technische Behauptung zu einer eigenen Technologie bleibt sprachlich unsicher.
\item Die Signatur wurde automatisch fehlerhaft erkannt.
\end{itemize}

\begin{asrbox}{Automatischer Pass 1 - Originaldatei}
Ja, okay und jetzt nochmal zurück zur Mengele des Denkens. Wenn wir also Informationen in die Zukunft schicken können und Informationen aus der Vergangenheit erhalten können, dann haben wir alle drei Zeiten in unserer Menge des Denkens drinnen. Dann sagt man auch, mein Hut da hat rein. dann sagt man auch mein Hut, der hat reihecken. Und zum Glück fliegt er nicht ständig weg der Hut. Jetzt muss man sich nochmal fragen, was ist denn Zeit? Und dann wissen wir, seit einstern Zeit und Raum sind gekoppelt. Und dann hatten die damals diese Idee mit diesem kausallicht kegel. Und das habe ich ja schon gesagt, dass es eine Vereinfachung ist, der Realität und allein die Unschärz Unschärferen Lationen für sich schon beweist, dass der Lichtkägel eine Vereinfachung der Realität ist, um es mal vorsichtig zu formulieren. Okay, wenn wir also jetzt bewiesenermaßen Informationen durch die Zeit schicken können, so wie wir gerade lustet darauf haben, dann kann ich meine Technologie meiner Technologie. Dann müssen wir uns noch mal fragen, was jenseits ist. Und wenn man dann mit der mänglere des Denkens arbeitet, dann stellt man fest, dass jenseits all das ist, was nicht hier und jetzt ist. Das heißt, wir können also mit dieser Technologie auch. mit dem Lähenseits kommunizieren. Quot errat demonstrantum Ingolph Lommern.
\end{asrbox}

\begin{asrbox}{Automatischer Pass 2 - normalisierter Cross-Check}
Ja, okay und jetzt noch mal zurück zur Mengele des Denkens. Wenn wir also Informationen in die Zukunft schicken können und Informationen aus der Vergangenheit erhalten können, dann haben wir alle drei Zeiten in unserer Menge des Denkens. drinnen, dann sagt man auch mein Hut der Hartei Ecken und zum Glück fliegt er nicht ständig weg der Hut. Jetzt muss man sich nochmal fragen, was ist denn Zeit und dann wissen wir seit einstern Zeit und Raum sind gekoppelt und dann hatten die damals die Und dann hatten die damals diese Idee mit diesem kausallicht Kegel. Und das habe ich ja schon gesagt, dass es eine Vereinfachung ist, der Realität und allein die unsherfere Lation für sich schon beweist, dass der Lichtkegel eine Vereinfachung der Realität ist, um es mal vorsichtig zu formulieren. Okay, wenn wir also jetzt bewiesen am Maßen Information jetzt bewiesenermaßen Informationen durch die Zeit schicken können. So wie wir gerade Lust darauf haben, dank meiner Technologie. Dann müssen wir uns noch mal fragen, was jenseits ist. Und wenn man dann mit der Mengele des Denkens arbeitet, dann stellt man fest, dass jenseits all das ist, was nicht hier uns jetzt ist. Das heißt, wir können also mit dieser Technologie auch, mit dem Jenseits kommunizieren. Quot Erwart. Quot errat demonstrandum engolf Lommann.
\end{asrbox}


\subsection{A06: Wessen ...}
\label{sec:a06}
\textbf{Konservativer semantischer Kern (kein Verbatim-Zitat):}
\begin{itemize}
\item Der Sprecher sagt, das Kraut gegen Dummheit habe einen Namen: „Steuer“.
\end{itemize}

\textbf{Wissenschaftliche Einordnung.} „Steuer“ kann im Deutschen Abgabe oder Steuerung bedeuten. Ohne Kontext, Zielgröße, Mechanismus und Evaluationsdesign folgt weder eine empirische Therapieaussage noch eine konkrete Steuerpolitik. Als Wortspiel verweist der Satz darauf, dass Verhalten steuerbar sein könnte.

\textbf{Status.} \status{APHORISTIC / lexically ambiguous / normative or control interpretation open}

\textbf{Akustische Unsicherheiten:}
\begin{itemize}
\item „einen Namen“ wird grammatisch fehlerhaft erkannt; der genaue gesprochene Kasus bleibt zu prüfen.
\item Die Signatur wurde automatisch fehlerhaft erkannt.
\end{itemize}

\begin{asrbox}{Automatischer Pass 1 - Originaldatei}
Das Kraut gegen Dummheit hat ein Namen. Steuer Quatt errat demonstrandum Ingolv Lohmann.
\end{asrbox}

\begin{asrbox}{Automatischer Pass 2 - normalisierter Cross-Check}
Das Kraut gegen Dummheit hat ein Namen. Steuer Quatt errat der Monsterandum Ingolv Lomann.
\end{asrbox}


\subsection{A07: Zu gut ...!}
\label{sec:a07}
\textbf{Konservativer semantischer Kern (kein Verbatim-Zitat):}
\begin{itemize}
\item Der Sprecher formuliert ein Wortspiel: Wer nicht „tutet“, solle „blasen“; sonst habe man keine Ahnung.
\end{itemize}

\textbf{Wissenschaftliche Einordnung.} Der Satz verbindet Tun, Blasen, Tuten und Ahnung in einer sprachlichen Pointe. Wissenschaftlich anschlussfähig wäre nur eine präzise Hypothese darüber, ob praktische Ausführung Kompetenz besser anzeigt als bloße Behauptung.

\textbf{Status.} \status{HUMOR / performance metaphor; no scientific result without operational definitions}

\textbf{Akustische Unsicherheiten:}
\begin{itemize}
\item Der zweite ASR-Pass verfehlt den ersten Halbsatz; die Erstlesung ist nur ein Kandidat.
\item Die Signatur wurde automatisch fehlerhaft erkannt.
\end{itemize}

\begin{asrbox}{Automatischer Pass 1 - Originaldatei}
Zu gut Deutsch. Wer nicht tutet, sollte blasen. Denn ansonsten hat man bewiesen, dass man gar keine Ahnung hat. Quart errat demonstrandum engolfloh man.
\end{asrbox}

\begin{asrbox}{Automatischer Pass 2 - normalisierter Cross-Check}
Zu gut Deutsch. Wenn ich tut, sollte Blasen. Denn ansonsten hat man bewiesen, dass man gar keine Ahnung hat. Quart errat demonstrandum engolfloh man.
\end{asrbox}

\section{Literatur und normative Quellen}
\begin{thebibliography}{99}
\bibitem{SHANNON-1948-I} C. E. Shannon, A Mathematical Theory of Communication, Bell System Technical Journal 27(3), 379-423 (1948). DOI: \href{https://doi.org/10.1002/j.1538-7305.1948.tb01338.x}{10.1002/j.1538-7305.1948.tb01338.x}.
\bibitem{SHANNON-1948-II} C. E. Shannon, A Mathematical Theory of Communication, Bell System Technical Journal 27(4), 623-656 (1948). DOI: \href{https://doi.org/10.1002/j.1538-7305.1948.tb00917.x}{10.1002/j.1538-7305.1948.tb00917.x}.
\bibitem{EINSTEIN-1905} A. Einstein, Zur Elektrodynamik bewegter Körper, Annalen der Physik 322(10), 891-921 (1905). DOI: \href{https://doi.org/10.1002/andp.19053221004}{10.1002/andp.19053221004}.
\bibitem{MINKOWSKI-1909} H. Minkowski, Raum und Zeit, Jahresbericht der Deutschen Mathematiker-Vereinigung 18, 75-88 (1909). \url{http://doi.library.cmu.edu/10.1184/OCLC/05366489}.
\bibitem{HEISENBERG-1927} W. Heisenberg, Über den anschaulichen Inhalt der quantentheoretischen Kinematik und Mechanik, Zeitschrift für Physik 43, 172-198 (1927). DOI: \href{https://doi.org/10.1007/BF01397280}{10.1007/BF01397280}.
\bibitem{NIST-CODATA-2022} NIST Standard Reference Database 121, 2022 CODATA recommended values of the fundamental physical constants, web version 9.0. \url{https://physics.nist.gov/constants}.
\bibitem{PEARL-2009} J. Pearl, Causality: Models, Reasoning, and Inference, 2nd ed., Cambridge University Press (2009).
\bibitem{EU-AI-ACT} Regulation (EU) 2024/1689 (Artificial Intelligence Act), Article 113 and applicable staged dates. \url{https://eur-lex.europa.eu/eli/reg/2024/1689/oj}.
\end{thebibliography}

\end{document}
